% Vereinfachte Architektur des modularen IMU-Sensorsystems (eingebunden in Expose.tex)
\begin{tikzpicture}[
    font=\small,
    >={Stealth[length=2mm]},
    comp/.style={draw=hkagray, fill=white, rounded corners=3pt, align=center,
                 minimum height=9mm, inner sep=3pt},
    layer/.style={draw=hkared, fill=hkared!12, rounded corners=3pt, align=center,
                  minimum width=5.2cm, minimum height=7mm},
    mod/.style={draw=hkagray, fill=hkagray!10, rounded corners=3pt, align=center,
                minimum width=2.0cm, minimum height=7mm, font=\footnotesize},
    part/.style={draw=hkagray, fill=white, rounded corners=2pt, align=center,
                 minimum width=1.5cm, minimum height=6mm, font=\footnotesize},
    group/.style={draw=hkared, line width=0.8pt, rounded corners=6pt, inner sep=7pt},
    newgroup/.style={draw=hkared, dashed, line width=0.8pt, rounded corners=6pt, inner sep=7pt},
    grouptitle/.style={fill=white, inner sep=2pt, font=\small\bfseries\color{hkared}},
    data/.style={->, draw=black!75},
    register/.style={->, dashed, draw=black!75},
]

% --- Sensorknoten ---
\foreach \name/\label/\y in {k1/Knoten 1/4.2, k2/Knoten 2/1.8, kn/Knoten N/-0.6} {
    \node[part] (\name imu) at (0,\y) {IMU};
    \node[part, minimum width=2.4cm] (\name fw) at (2.6,\y) {Zephyr-Firmware\\Sensor-API};
    \node[part] (\name ble) at (5.0,\y) {BLE};
    \draw[->, draw=black!75] (\name imu) -- (\name fw);
    \draw[->, draw=black!75] (\name fw) -- (\name ble);
}
\begin{scope}[on background layer]
    \node[group, fit=(k1imu)(k1ble)] (g1) {};
    \node[group, fit=(k2imu)(k2ble)] (g2) {};
    \node[newgroup, fit=(knimu)(knble)] (gn) {};
\end{scope}
\node[grouptitle] at (g1.north) {Knoten 1 (Oberarm)};
\node[grouptitle] at (g2.north) {Knoten 2 (Unterarm)};
\node[grouptitle] at (gn.north) {Knoten N (neu)};
\node[font=\Large, text=hkagray] at (2.6,0.65) {$\vdots$};

% --- Zentraleinheit ---
\node[layer] (trans) at (10.5,4.2) {\textbf{Transport} (BLE)};
\node[layer] (reg)   at (10.5,3.0) {\textbf{Register}\\[-1pt]\footnotesize Selbstbeschreibungen};
\node[layer] (sync)  at (10.5,1.8) {\textbf{Synchronisation}\\[-1pt]\footnotesize gemeinsame Zeitachse};
\node[mod] (m1) at (8.3,0.4) {Orientierung};
\node[mod] (m2) at (10.5,0.4) {Visualisierung};
\node[mod] (m3) at (12.7,0.4) {Aufzeichnung};
\node[mod, dashed, draw=hkared] (m4) at (10.5,-0.6) {weiteres Modul};
\node[comp, minimum width=2.0cm] (cfg) at (15.6,3.0) {\textbf{Konfiguration}\\[1pt]\footnotesize Knoten $\to$ Segment};

\draw[data] (trans) -- (reg);
\draw[data] (reg) -- (sync);
\draw[data] (sync.south) -- ++(0,-0.35) -| (m1.north);
\draw[data] (sync.south) -- (m2.north);
\draw[data] (sync.south) -- ++(0,-0.35) -| (m3.north);
\draw[->, draw=hkagray] (cfg.west) -- (reg.east);

\begin{scope}[on background layer]
    \node[group, fit=(trans)(m1)(m3)(m4)] (gz) {};
\end{scope}
\node[grouptitle] at (gz.north) {Zentraleinheit (PC/Smartphone)};

% --- Verbindungen Knoten -> Zentraleinheit ---
\draw (k1ble.east) -- (6.6,4.2);
\draw (k2ble.east) -- (6.6,1.8);
\draw[dashed] (knble.east) -- (6.6,-0.6);
\draw (6.6,4.2) -- (6.6,-0.6);
\draw[data] (6.6,4.2) -- (trans.west);
\node[font=\footnotesize, align=center, text=hkagray, fill=white, inner sep=1pt] at (6.6,2.9)
    {Anmeldung\\+ Daten};

\end{tikzpicture}
